%% Extraído de docs/TCC_documento/Textuais/Metodologia/02_MotorBLDC.tex (linhas 15-50)
    \begin{tikzpicture}[transform shape, american]
        % Ponto Neutro
        \node[circ] (N) at (0,0) {};
        \node[above left, xshift=-0.1cm] at (N) {$N$};

        % Fase U (Para Cima)
        \draw (N) to[american resistor, l=$R_{fase}$] (0,2)
                  to[cute inductor, l=$L_{fase}$] (0,4)
                  node[ocirc] (U) {}
                  node[above] {Fase $U$};

        % Fase V (Para a Direita/Baixo) - Ângulo de -30°
        \draw (N) to[american resistor, l=$R_{fase}$] (1.732,-1)
                  to[cute inductor, l=$L_{fase}$] (3.464,-2)
                  node[ocirc] (V) {}
                  node[right] {Fase $V$};

        % Fase W (Para a Esquerda/Baixo) - Ângulo de 210°
        \draw (N) to[american resistor, l_=$R_{fase}$] (-1.732,-1)
                  to[cute inductor, l_=$L_{fase}$] (-3.464,-2)
                  node[ocirc] (W) {}
                  node[left] {Fase $W$};

        % --- Indicação do Ensaio (Medição Linha a Linha) ---

        % Linhas de chamada (tracejadas) saindo de U e V para a direita
        \draw[dashed, blue!70!black] (U) -- (5.5, 4);
        \draw[dashed, blue!70!black] (V) -- (5.5, -2);

        % Seta dupla indicando a medição com texto explicativo
        \draw[<->, blue!70!black, thick] (5.5, 3.9) -- node[fill=white, align=center, text=black] {Ensaio Linha a Linha \\ $V_{DC}, I_{DC}$ (Resistência) \\ $L_{linha}$ (Indutância) \\ \small \textit{Equivalente a 2 em série}} (5.5, -1.9);

        % Marcação do caminho da corrente para ficar mais didático (Opcional)
        \draw[->, red, thick, rounded corners=5mm] (0.5, 3.5) -- (0.5, 0.5) -- (3.0, -1.0) node[pos=0.8, above right] {$I_{ensaio}$};

    \end{tikzpicture}
